\documentclass[12pt,a4paper]{article}
\usepackage[margin=25mm, headheight=15pt, headsep=10pt]{geometry}
\usepackage{fontspec}
\usepackage[table]{xcolor} % Note: table option is required for rowcolors
\usepackage{titlesec}
\usepackage{tocloft}
\usepackage{fancyhdr}
\usepackage{booktabs}
\usepackage{tcolorbox}
\tcbuselibrary{skins, breakable}
\usepackage[colorlinks=true, linkcolor=emerald, urlcolor=emerald, bookmarks=true]{hyperref}

\definecolor{emerald}{RGB}{0,102,51}
\definecolor{darkred}{RGB}{139,0,0}
\definecolor{lightgray}{RGB}{245,245,245}

\usepackage[nil]{babel}
\babelprovide[import, main]{bengali}
\babelprovide[import]{arabic}
\babelfont{rm}[Renderer=HarfBuzz,Scale=1.0]{Noto Serif Bengali}
\babelfont{sf}[Renderer=HarfBuzz]{Noto Sans Bengali}
\babelfont{tt}[Renderer=HarfBuzz]{Noto Sans Bengali}
\babelfont[arabic]{rm}[Renderer=HarfBuzz,Scale=1.1]{Noto Naskh Arabic}

% Section styling
\titleformat{\section}{\Large\bfseries\color{emerald}}{\thesection.}{0.5em}{}
\titleformat{\subsection}{\large\bfseries\color{emerald!85!black}}{\thesubsection.}{0.5em}{}
\titleformat{\subsubsection}{\normalsize\bfseries\color{emerald!70!black}}{\thesubsubsection.}{0.5em}{}
\titlespacing*{\section}{0pt}{18pt}{8pt}

% Fancy Header/Footer setup
\pagestyle{fancy}
\fancyhf{} % clear all header and footer fields
\fancyhead[L]{\color{emerald}\small\textbf{\nouppercase{\leftmark}}}
\fancyfoot[L]{\scriptsize\color{black!50}{\lat{AI}}-সংকলিত, আলেম-যাচাইকৃত নয়}
\fancyfoot[C]{\color{emerald!70!black}\thepage}
\renewcommand{\headrulewidth}{0.4pt}
\renewcommand{\headrule}{\hbox to\headwidth{\color{emerald}\leaders\hrule height \headrulewidth\hfill}}
\renewcommand{\footrulewidth}{0pt}

% Custom boxes
% 1. Ayat/Hadith Box
\newtcolorbox{ayatbox}[1][]{
    enhanced, breakable, 
    colback=emerald!5, 
    colframe=emerald, 
    boxrule=0.5pt, 
    arc=3pt, 
    left=10pt, right=10pt, top=10pt, bottom=10pt,
    #1
}

% 2. Quote/Translation Box
\newtcolorbox{quotebox}[1][]{
    enhanced, breakable,
    frame hidden,
    colback=lightgray,
    borderline west={3pt}{0pt}{emerald},
    left=15pt, right=10pt, top=10pt, bottom=10pt,
    sharp corners,
    #1
}

% 3. AI-generated notice box
\newtcolorbox{warnbox}[1][]{
    enhanced, breakable,
    colback=red!4,
    colframe=darkred,
    boxrule=0.8pt,
    arc=3pt,
    left=10pt, right=10pt, top=10pt, bottom=10pt,
    #1
}

% Commands
\newcommand{\arab}[1]{\foreignlanguage{arabic}{#1}}
\newcommand{\kib}[1]{\textcolor{emerald}{\textbf{#1}}}

% Latin (English) fallback: Noto Serif Bengali has NO Latin glyphs
\newfontfamily\latfont[Renderer=HarfBuzz]{Noto Serif}
\newcommand{\lat}[1]{{\latfont #1}}

\setlength{\parskip}{5pt}
\setlength{\parindent}{0pt}

% Fix for missing bullet in Noto Serif Bengali
\renewcommand{\labelitemi}{$\bullet$}



\begin{document}
\begin{center}{\Large\bfseries\color{emerald} পার্ট ৫ক — মাসআলা: জামাত, ইমাম ও মুক্তাদি, মাসবুক}\end{center}
\vspace{1em}
\section*{পার্ট ৫ক — মাসআলা: জামাত, ইমাম ও মুক্তাদি, মাসবুক}

\begin{warnbox}
 \textbf{গুরুত্বপূর্ণ নোটিশ (\lat{Important} \lat{Notice}):} এই কন্টেন্টটি \lat{AI} (কৃত্রিম বুদ্ধিমত্তা) দিয়ে প্রস্তুত ও সংকলিত। \lat{AI} কঠোর সূত্র-মান নীতি মেনে শুধু নির্ভরযোগ্য উৎস থেকে তথ্য সংগ্রহ করেছে — কেবল সহীহ/হাসান হাদিস ও সহীহ সনদের আসার রাখা হয়েছে, দুর্বল সনদ বাদ দেওয়া হয়েছে। তবে এটি কোনো আলেম বা মুহাদ্দিস কর্তৃক যাচাইকৃত নয়।
\end{warnbox}

\begin{quote}
\textbf{সম্পর্কিত ফাইল:} পার্ট ২ (পূর্বশর্ত), পার্ট ৩ (শ্রেণিবিভাগ), পার্ট ৫খ (সাহু সিজদা), পার্ট ৫গ (সফর-অসুস্থ-নারী), পার্ট ৫ঘ (জুমা-ঈদ), পার্ট ৬ (মাযহাব-তুলনা)।
\end{quote}

\begin{quote}
\textbf{যে প্রশ্নের উত্তর এই অংশে:} জামাতের গুরুত্ব কত? কখন জামাত ওয়াজিব? ইমামের শর্ত কী? মুক্তাদি ইমামের পেছনে কীভাবে চলবে? দেরিতে (মাসবুক) পৌঁছালে কী করতে হবে? নারীর কি জামাত-ফজিলত আছে?
\end{quote}

\medskip\hrule\medskip

\subsection*{১. জামাতের ফজিলত ও গুরুত্ব}

\textbf{সহীহ হাদিস (সহীহ বুখারি ৬৪৫; সহীহ মুসলিম ৬৫০):}

\begin{quote}
আবু হুরাইরা (রাঃ) থেকে: নবীজি (সা.) বলেন — \textbf{\lat{“}জামাতের নামাজ একাকী নামাজের চেয়ে পঁচিশ বা সাতাশ গুণ উত্তম।\lat{”}}
\end{quote}

\textbf{সহীহ হাদিস (সহীহ মুসলিম ৬৫৩):}

\begin{quote}
আবদুল্লাহ ইবনে উম্মে মাকতুম (রাঃ) — যিনি অন্ধ ছিলেন — নবীজি (সা.)-কে বললেন: হে আল্লাহর রাসূল! আমার ঘরে মসজিদ থেকে দূরে; আমি চাই ঘরেই নামাজ পড়ি। নবীজি (সা.) বললেন: \textbf{\lat{“}আপনি কি আযান শুনতে পান?\lat{”}} বললেন: হ্যাঁ। নবীজি (সা.) বললেন: \textbf{\lat{“}তবে উত্তর দাও (মসজিদে এসো)।\lat{”}}
\end{quote}

\begin{quote}
\textbf{ব্যাখ্যা:} অন্ধ সাহাবি থেকেও জামাত ছাড়ার অনুমতি নিলেন না — দূরত্ব ও অক্ষমতা সত্ত্বেও। এ থেকেই ফকিহরা জামাতের \textbf{কঠোর গুরুত্ব} বুঝেছেন (নিচে স্তরভেদ)।
\end{quote}

\textbf{সহীহ হাদিস (সহীহ বুখারি ৬৪৪; সহীহ মুসলিম ৬৫১):}

\begin{quote}
আবু হুরাইরা (রাঃ) থেকে: নবীজি (সা.) বলেন — \textbf{\lat{“}আমার মন চায় লোকদের জামাতের জন্য নির্দেশ দিই, তারপর ঘোষণা দিই: যারা (বিনা ওজরে) ঘরে বসে জামাত ত্যাগ করে... (তাদের ঘরে আগুন জ্বালিয়ে দেওয়ার কথা বললেন)।\lat{”}} (এটি জামাত ত্যাগের ভয়াবহ সতর্কতা)
\end{quote}

\subsubsection*{১.১ জামাতের হুকুম — মাযহাবগত অবস্থান}

\begin{table}[h]\centering\small
\begin{tabular}{ll}
মাযহাব & ফরজ নামাজে জামাতের হুকুম \\
\hline
\textbf{হাম্বলি} & ফরজে আইন (পুরুষদের জন্য আবশ্যক) — প্রসিদ্ধ মত \\
\textbf{হানাফি} & ওয়াজিব (প্রায় ফরজের মতো; সুন্নতে মুআক্কাদা এক মত) \\
\textbf{শাফেয়ি} & সুন্নতে মুআক্কাদা (জামাতের জন্য জোরালো তাগিদ; ফরজ কেফায়া এক মত) \\
\textbf{মালেকি} & সুন্নতে মুআক্কাদা (কিছু বর্ণনায় ওয়াজিব) \\
\hline
\end{tabular}
\end{table}

\begin{quote}
\textbf{সারমর্ম:} সব মাযহাবের মিল — পুরুষ মুকিমের জন্য ফরজ নামাজ \textbf{মসজিদে জামাতে} পড়া অত্যন্ত গুরুত্বপূর্ণ; শিথিলতার স্তরই আলাদা। \textbf{ব্যতিক্রম (সব মাযহাবে):} রোগ, প্রবল বৃষ্টি, ভয়, ঘুমের চাপ, ওজু-ভঙ্গের আশঙ্কা, অসহ্য জরুরত ইত্যাদি বৈধ ওজরে জামাতে যাওয়া মাফ।
\end{quote}

\medskip\hrule\medskip

\subsection*{২. ইমাম — কে, কাকে, কীভাবে}

\subsubsection*{২.১ ইমাম হওয়ার শর্ত}

\begin{enumerate}
\item মুসলিম, বালেগ, সুস্থ বুদ্ধি — মৌলিক শর্ত (সব মাযহাবে)।
\item \textbf{পুরুষের ইমামতি নারীদের জন্য} — নারী ইমামতি পুরুষের জন্য বৈধ নয় (জমহুরের সর্বসম্মত অবস্থান; নারী নারীদের ইমামতি — পার্ট ৫গ)।
\item পবিত্র (ওযু/গোসল/তায়াম্মুম) — শর্ত (পার্ট ২)।
\item \textbf{সঠিকভাবে কুরআন তিলাওয়াত} করতে পারেন — ইমামের জন্য অপরিহার্য; তিলাওয়াতের ত্রুটি নামাজের কিরাআত নষ্ট করলে ইমামতি সমস্যা (বিস্তারিত পার্ট ৬)।
\item \textbf{অগ্রাধিকার (ইমাম বাছাইয়ের আদব):} নবীজি (সা.) বলেছেন — \textbf{\lat{“}লোকদের ইমাম করো তাদের কুরআন-বিশেষজ্ঞকে\lat{”}} (সহীহ মুসলিম ৬৭৩)। পর্যায়ক্রমে — বেশি জ্ঞান, বেশি ধর্মপরায়ণ, বড়, সুন্দর-শুদ্ধ তিলাওয়াত।
\end{enumerate}
\subsubsection*{২.২ ইমামের কর্তব্য}

\textbf{সহীহ হাদিস (সহীহ মুসলিম ৬৭৫):}

\begin{quote}
আবু হুরাইরা (রাঃ) থেকে: নবীজি (সা.) বলেন — \textbf{\lat{“}যদি তোমাদের কেউ লোকদের ইমামতি করে, তবে সে যেন (নামাজ) সংক্ষিপ্ত (উত্তমভাবে) করে; কেননা তাদের মধ্যে দুর্বল, অসুস্থ ও বৃদ্ধও থাকে। আর একাকী নামাজ পড়লে যত খুশি দীর্ঘ করতে পারে।\lat{”}}
\end{quote}

\textbf{সহীহ হাদিস (সহীহ বুখারি ৭২২; সহীহ মুসলিম ৪২৩):}

\begin{quote}
নবীজি (সা.) বলেন — \textbf{\lat{“}ইমাম তো এই জন্যই বানানো হয়েছে যে তাকে অনুসরণ করা হবে।\lat{”}} — এরপর কাতার সোজা করার, রুকু-সিজদা ইমামের সাথে সাথে করার বিস্তারিত নির্দেশ।
\end{quote}

\subsubsection*{২.৩ ইমামের আগে-পরে চলা — কঠোর সতর্কতা}

\textbf{সহীহ হাদিস (সহীহ বুখারি ৬৯১; সহীহ মুসলিম ৪২৭):}

\begin{quote}
আবু হুরাইরা (রাঃ) থেকে: নবীজি (সা.) বলেন — \textbf{\lat{“}তোমাদের কেউ কি ভয় করে না যে, সে ইমামের আগে মাথা তুললে আল্লাহ তার মাথা গাধার মাথা বানিয়ে দেবেন?\lat{”}} — ইমামের আগে রুকু/সিজদা করা \textbf{হারাম-সদৃশ গুরুতর}।
\end{quote}

\medskip\hrule\medskip

\subsection*{৩. মুক্তাদির নিয়ম — ইমামের পেছনে কী করবেন}

\subsubsection*{৩.১ সাধারণ নিয়ম}

\begin{enumerate}
\item \textbf{তাকবিরে তাহরিমা} মুক্তাদি ইমামের পরপরই দেবেন (ইমামের পরে); ইমামের আগে নয়।
\item \textbf{রুকু-সিজদা-সালাম} — সব মুক্তাদি ইমামের \textbf{পরে} বা \textbf{সাথে সাথে}; কখনো আগে নয়।
\item \textbf{মুক্তাদির কিরাআত:}
\end{enumerate}
\begin{itemize}
\item \textbf{জোরে (জাহরি) নামাজে} (ফজর, মাগরিব, ইশা): মুক্তাদি চুপ থাকে, ইমামের কিরাআতই যথেষ্ট।
\item \textbf{নিরবে (সিররি) নামাজে} (যোহর, আসর): হানাফি মতে মুক্তাদি চুপ থাকে; জমহুর (শাফেয়ি-মালেকি) মতে ফাতিহা পড়া উচিত — হাম্বলি মতেও ফাতিহা পড়া মুস্তাহাব। (এই ইখতিলাফটি বাস্তব; ইমাম-মুক্তাদি দুই পক্ষই বৈধ — দলিল দুই রকম, পার্ট ৬-এ বিশ্লেষণ)
\end{itemize}
\begin{enumerate}
\item \textbf{সালাম} — ইমামের পর।
\item \textbf{মাসবুক হলে} — নিচে ৪ নম্বর অংশ।
\end{enumerate}
\subsubsection*{৩.২ মুক্তাদি পিছনে ছুটে গেলে (ইমাম এগিয়ে গেলে)}

\begin{itemize}
\item রুকু/সিজদা ইমামের \textbf{আগে} শুরু করা হারাম-সদৃশ (বুখারি ৬৯১)।
\item তবে \textbf{ইমামের পরেই} রুকু/সিজদা করলে নামাজ বৈধ — এমনকি ইমামের পরে দ্বিতীয় সিজদা মিস হলে মাযহাবগত সমাধান আছে (পার্ট ৫খ)।
\end{itemize}
\medskip\hrule\medskip

\subsection*{৪. মাসবুক (দেরিতে জামাতে পৌঁছা) — পূর্ণ পদ্ধতি}

\textbf{মাসবুক} = যে মুক্তাদি প্রথম রাকাতের রুকু পেরিয়ে জামাতে ঢুকেছে।

\textbf{সহীহ হাদিস (সহীহ বুখারি ৬৩৫; সহীহ মুসলিম ৬০২):}

\begin{quote}
আবু হুরাইরা (রাঃ) থেকে: নবীজি (সা.) বলেন — \textbf{\lat{“}...যা (রাকাত) জামাত থেকে পেলে তা পড়ো, আর যা মিস হলে তা পূর্ণ করো।\lat{”}}
\end{quote}

\subsubsection*{৪.১ মাসবুকের মূল নিয়ম (সব মাযহাবের মোটা দাঁড়া)}

\begin{itemize}
\item জামাতে যে রাকাত \textbf{ইমামের সাথে (রুকুসহ) পেয়েছেন} — সেগুলো তার "জামাতের রাকাত"।
\item যে রাকাত মিস হয়েছে — ইমাম সালাম ফিরানোর পর \textbf{নিজে পূর্ণ করবেন} (বাকি রাকাত একাকী-নিয়মে)।
\item \textbf{নিয়ত ও প্রথম রাকাত:} মাসবুকের প্রথম তাকবিরেই তার নিজের প্রথম রাকাত — জামাতে সে যত রাকাত পায়, তা তার শেষের রাকাত।
\end{itemize}
\subsubsection*{৪.২ ধাপে ধাপে উদাহরণ (৪ রাকাতের নামাজ, ২ রাকাত পাওয়া)}

\begin{enumerate}
\item ইমামের সাথে তাকবিরে তাহরিমা $\rightarrow$ ইমামের সাথে ২য় রাকাতে (রুকুসহ) পেলেন। এই ২ রাকাত = জামাতের রাকাত।
\item ইমাম সালাম ফিরালে \textbf{মাসবুক বসে থাকবেন না} — উঠে \textbf{আরও ২ রাকাত} নিজে পড়বেন (অবশিষ্ট ফরজ পূর্ণ করতে)।
\item বাকি ২ রাকাতে: প্রথমটিতে ফাতিহা+সূরা, দ্বিতীয়টিতে ফাতিহা+সূরা (নিজের নিয়মে) $\rightarrow$ শেষে তাশাহহুদ-সালাম।
\end{enumerate}
\begin{quote}
\textbf{সূক্ষ্ম বিষয়:} মাসবুকের \textbf{দ্বিতীয় জামাত-রাকাত} ইমামের সাথে; তার \textbf{বাড়তি রাকাতগুলো} সে ইমামের পরে একাকী পড়ে — এতে সে \textbf{কোনো রাকাতে যেন তাশাহহুদ/সালাম ভুল না করে} (সাধারণ মাসবুকের ভুল — পার্ট ৭খ)।
\end{quote}

\subsubsection*{৪.৩ মাসবুক ইমামের ২য় রাকাতে পেল — কীভাবে মিলবে?}

\begin{itemize}
\item \textbf{হানাফি:} মাসবুক ইমামের ২য় রাকাতে তাশাহহুদে বসে ইমামের সাথে। ইমাম সালামের পর ৩ ও ৪ রাকাত নিজে পড়ে — \textbf{কিন্তু প্রথম রাকাত নিজের} (ফাতিহা+সূরা)।
\item \textbf{জমহুর:} একই মূলনীতি — যে রাকাত পেল, সেটাই তার শেষ রাকাত; বাকি নিজে।
\end{itemize}
\begin{quote}
\textbf{মাসবুকের গুরুত্বপূর্ণ সতর্কতা:} ইমামের পেছনে থাকা অবস্থায় \textbf{মাসবুক কখনো নিজের "বাকি রাকাত" শুরু করবেন না} — বরং ইমামের সাথে সম্পূর্ণ অনুসরণ, তারপর পূর্ণতা। এটি ভুল করলে নামাজের গঠন নষ্ট হয়ে যায় (পার্ট ৭খ-এর ভুল-তালিকায় আছে)।
\end{quote}

\medskip\hrule\medskip

\subsection*{৫. কাতার (সারি) — কাতার সোজা করার নির্দেশ}

\textbf{সহীহ হাদিস (সহীহ বুখারি ৭২৩; সহীহ মুসলিম ৪৩২):}

\begin{quote}
আনাস (রাঃ) থেকে: নবীজি (সা.) বলেছেন — \textbf{\lat{“}তোমাদের কাতার সোজা করো, কেননা কাতার সোজা করাই নামাজের সৌন্দর্যের অংশ।\lat{”}}
\end{quote}

\begin{itemize}
\item কাতার সোজা করা, কাঁধে কাঁধ মিলানো ও ফাঁকা জায়গা পূরণ করা \textbf{সুন্নত} (এবং এক মতে ওয়াজিব)।
\item \textbf{প্রথম কাতারের ফজিলত:} \lat{“}প্রথম কাতারে আল্লাহ ও ফেরেশতারা দরুদ পড়েন\lat{”} (সুনানে আবু দাউদ ৬৬০ — সহীহ)। নবীজি (সা.) প্রথম কাতারে প্রার্থনা করেছেন — (সহীহ বুখারি ৬১৫, সহীহ মুসলিম ৪৩৭)।
\end{itemize}
\medskip\hrule\medskip

\subsection*{৬. নারী ও জামাত}

\textbf{সহীহ হাদিস (সহীহ বুখারি ৯০০; সহীহ মুসলিম ৮৮৪):}

\begin{quote}
নবীজি (সা.) বলেছেন — \textbf{\lat{“}নারীদের মসজিদে যেতে বাধা দিও না।\lat{”}}
\end{quote}

\begin{itemize}
\item নারীর মসজিদে (পেছনের কাতারে, পর্দা রক্ষা করে) নামাজ জায়েয ও ফজিলতপূর্ণ — নবীজির যুগের আমল থেকে প্রমাণিত (উম্মে আতিয়্যাহর ঘটনা — পার্ট ৮-এর আলাদা ফাইল)।
\item তবে \textbf{একাকী নফল বাড়ি} পড়া নারীর জন্য উত্তম বলে একটি হাদিস আছে — \lat{“}নারীর নামাজ তার ঘরে তার উঠানের চেয়ে উত্তম...\lat{”} (সহীহ বুখারি ৮৭১; সহীহ মুসলিম ১০৩৫) — এটি \textbf{নফল/বাড়তি} নামাজের জন্য; \textbf{ফরজ জামাতে} নারীর অংশগ্রহণ জায়েয, বিশেষত নবীজি-যুগের ঈদ ও জুমায় নারীদের উপস্থিতি প্রমাণিত।
\item \textbf{বিস্তারিত:} পার্ট ৫গ-এ।
\end{itemize}
\medskip\hrule\medskip

\subsection*{৭. প্রশ্নোত্তর (এক নজরে)}

\begin{table}[h]\centering\small
\begin{tabular}{ll}
প্রশ্ন & উত্তর \\
\hline
জামাত কি ওয়াজিব? & হাম্বলি-হানাফি: ওয়াজিব/ফরজ; শাফেয়ি-মালেকি: সুন্নতে মুআক্কাদা — সবাই একমত: অযৌক্তিক ত্যাগ বড় গুনাহ \\
ইমাম কে হবে? & পবিত্র, তিলাওয়াতে সুদক্ষ, ধার্মিক, বয়সে বড় — এই ক্রমে (বুখারি-মুসলিমের রূপরেখায়) \\
ইমামের আগে চলা & হারাম-সদৃশ — কঠোর সতর্কতা (বুখারি ৬৯১) \\
মাসবুক কী করবে? & যা পেল = জামাতের রাকাত; যা মিস = ইমামের সালামের পর নিজে পূর্ণ করবে (বুখারি ৬৩৫) \\
নারীর মসজিদে যাওয়া & জায়েয ও ফজিলতপূর্ণ (বুখারি ৯০০) \\
\hline
\end{tabular}
\end{table}

\medskip\hrule\medskip

\subsection*{৮. রেফারেন্স}

\begin{itemize}
\item সহীহ বুখারি — ৬১৫, ৬৩৫, ৬৪৪, ৬৪৫, ৬৯১, ৭২২, ৭২৩, ৮৭১, ৯০০
\item সহীহ মুসলিম — ৬০২, ৬৫০, ৬৫১, ৬৫৩, ৬৭৩, ৬৭৫, ৪২৩, ৪২৭, ৪৩২, ৮৮৪, ১০৩৫
\item সুনানে আবু দাউদ — ৬৬০ (প্রথম কাতারের দরুদ)
\item আল-ফিকহুল ইসলামি ওয়া আদিল্লাতুহু (আল-যুহাইলি) — মাযহাব-অবস্থান সারণি
\end{itemize}
\begin{quote}
\textbf{দ্রষ্টব্য:} মাসবুকের জটিলতম ক্ষেত্র (যেমন — ৩য় রাকাতে ঢোকা, ইমাম সাহু সিজদা দিলে মাসবুকের কী, সাহু সিজদায় মাসবুকের অংশ) পার্ট ৫খ-এ; ইমাম-মুক্তাদির কিরাআতের মাযহাব-তুলনা পার্ট ৬-এ; নারী-জামাতের বিশদ পার্ট ৫গ-এ।
\end{quote}

\end{document}
